\section{Introduction}

Predicting the direction of Earth from a lunar station connects several familiar
astronomical operations: evaluating an ephemeris, identifying the coordinate
origin, rotating into the Moon's body-fixed frame and projecting a finite line
of sight onto a local horizon. The calculation also connects distinct kinds of
scientific information. Its epoch has a time scale; its vector has units, a
frame and a center; and its observer belongs to a reference surface with a
particular orientation model. Preserving these relationships is a central
software design problem in astronomical computing.

Scientific libraries address this problem through quantity objects, coordinate
classes, frame metadata and domain-specific interfaces. As workflows grow,
these representations must also compose across numerical components and
application boundaries. A useful architecture should let a developer reuse a
coordinate container or time representation while retaining the information
needed by the next scientific operation. This supports both the development of
new algorithms and the assembly of reliable astronomical applications.

We present Siderust~\cite{Valles_Puig_Siderust_Type-Safe_Astronomy}, an open-source Rust ecosystem organized around this form
of scientific composition. We separate reusable quantity, geometry, time and
numerical abstractions from astronomy-specific models and workflows. The domain
crate combines these components into reference-system transformations,
ephemerides, surface observations and astrodynamical calculations. Scientific
concepts are visible in public types, while explicit contexts supply values
such as observer parameters, Earth orientation and ephemeris choices.

Rust provides a practical basis for this design. Generic types express units,
time scales, frames and centers; traits define compatible scientific
operations; and ownership and borrowing support controlled sharing of data.
Consequently, a frame transformation can change a vector's representation
explicitly, a time conversion can preserve its scale information, and a surface
observation can retain its celestial-body identity. These mechanisms make
assumptions visible to both application authors and the compiler.

Siderust introduces a composable scientific-software architecture in which
reusable mathematical and numerical components are connected through explicit
astronomical semantics. Our contribution has three complementary parts.
First, we design and implement
a modular scientific architecture in which independently reusable mathematical
components are composed through a common astronomical domain layer. Second,
we provide expressive interfaces that carry physical and astronomical semantics
across component boundaries and permit static checking of selected constraints.
Third, we demonstrate integrated terrestrial and lunar observing workflows and
evaluate representative transformations, ephemerides and computational costs
against established references and matched arithmetic baselines.

These parts form one design: reusable components supply the operations, semantic
interfaces preserve their interpretation, and integrated applications test their
composition. The scientific contribution is this implemented connection across
layers, rather than a new ephemeris or a claim that typed quantities are novel.
Researchers can build new domain calculations on the same representations
without reimplementing quantities, time conversion or local geometry. Section~2 positions the design among
related systems. Sections~3--5 describe composition, public interfaces and
language integration. Section~6 develops the observing applications;
Sections~7--8 present the evaluation and results. We discuss architectural
implications and future directions in Section~9.
